\section{Connection to work on data augmentations}

Recently, there have been two papers published on using data augmentations for reinforcement learning, RAD~\cite{laskin2020reinforcement} and DrQ~\cite{kostrikov2020image}. These two papers present the version of CURL without an auxiliary contrastive loss but rather directly feeding in the augmented views of the image observations to the underlying value / policy network(s). Both RAD and DrQ present results on both continuous and discrete control environments, surpassing the results presented in CURL on both the DMControl and Atari benchmarks. Plenty of researchers have opined in public forums whether the results in RAD and DrQ make CURL irrelevant if the objective is to use data augmentations for data-efficient reinforcement learning. We believe that answering this question needs more nuance and present our opinions below:

1. If one has access to a rich stream of rewards from the underlying environment and is interested in optimzing the performance in terms of average reward, RAD and DrQ are likely to work better than CURL. The reason for this is simply that RAD and DrQ directly optimize for the objective one cares about, while CURL introduces an additional auxiliary consistency objective.

2. If one {\it does not} have access to a rich stream of rewards and is interested in learning good latent spaces in a {\it task agnostic manner} that can allow for data-efficient controllers across multiple tasks, CURL is the only option since the contrastive objective in CURL is reward independent. Our ablation on the detached encoder with the CURL objective present evidence that one could build simple MLPs on top of the CURL features without fine-tuning the underlying encoder and still be data-efficient on many of the DMControl tasks.

3. Future work in data-efficient reinforcement learning, particularly for real world settings, is likely to require approaches that {\it do not rely on reward functions}. In such scenarios, CURL is likely to be the more preferred approach. Further, one could potentially use CURL in a scenario where unsupervised pre-training without reward functions is initially performed before fine-tuning to the RL objective across multiple tasks.

Given the above reasons, there isn't a straightforward answer as to which is the better algorithm and the answer really depends on what the researcher / practioner wants to solve. We also emphasize that CURL was the first approach that used data augmentations effectively to significantly improve the data-efficiency of model-free reinforcement learning methods with very simple changes and showed improvement over relatively more complex model-based methods. The augmentations and results in CURL inspired future work in the form of RAD and DrQ. We hope that the analysis and results presented in CURL encourage researchers to employ data augmentations, contrastive losses and unsupervised pre-training for future reinforcement learning research.